\chapter{Route B: Ambient Solid Six Operations and Constructible Descent}
\label{chap:proetale-six}

Let \(P\) be a set of primes and put

\[
  \widehat{\mathbf Z}_{P}:=\prod_{r\in P}\mathbf Z_r.
\]

For the Weil conjectures we take \(P=\{\ell\}\), so that
\(\widehat{\mathbf Z}_{P}=\Zl\), and every scheme is a qcqs
\(\mathbf Z[1/\ell]\)-scheme.  This chapter adopts \emph{Route B}:

\[
\boxed{
 \begin{gathered}
 \text{one presentable ambient category of pro-etale sheaves}\\
 \Downarrow\\[-0.4em]
 \text{a coherent six-functor formalism}\\
 \Downarrow\\[-0.4em]
 \text{restriction to the constructible subcategory}.
 \end{gathered}}
\tag{4.B}\label{eq:route-b-overview}
\]

No operation is \emph{defined} by first constructing a compatible family
modulo \(\ell^n\) and then taking its inverse limit.  There is nevertheless
an essential correction to the phrase ``all pro-etale sheaves'': the naive
unrestricted derived category is too large.  The correct ambient category is
the reflective localization of \emph{solid} pro-etale complexes.  Solidity
is an ambient continuity condition, not a constructibility or finite
stratification condition.

\section{The unrestricted category is impossible}

Write
\[
 \Dnaive(X,\Lambda)
   :=\D\bigl(\operatorname{Sh}(X_{\proet},\Lambda)\bigr)
\]
for the unbounded derived category of all pro-etale \(\Lambda\)-module
sheaves (or its hypersheaf enhancement).

\begin{theorem}[Failure of unrestricted proper base change]
\label{thm:naive-proetale-pbc-fails}
Let \(r\) be a prime.  In the cartesian square
\begin{equation}
\label{eq:naive-pbc-counterexample}
\begin{tikzcd}
 \PP^1_{\mathbf C} \arrow[r,"q"] \arrow[d,"g"'] &
 \PP^1_{\mathbf Q} \arrow[d,"f"] \\
 \Spec(\mathbf C) \arrow[r,"p"'] & \Spec(\mathbf Q),
\end{tikzcd}
\end{equation}
the proper base-change transformation
\[
 p^*f_*\longrightarrow g_*q^*
\]
is not an equivalence on
\(\Dnaive(\PP^1_{\mathbf Q},\mathbf F_r)\).
\end{theorem}

\begin{proof}
This is \cite[Remark~7.15]{HaineEtAlCondensedHomotopy}.  We recall the
contradiction because it identifies exactly what goes wrong.  If the displayed
map were an equivalence, then its mate under the left adjoints would be an
equivalence.  Since pullback from \(\Spec(\mathbf Q)\) to
\(\Spec(\mathbf C)\) has the expected left adjoint on these pro-etale
hypersheaf categories, applying the mate to the unit and then taking relative
homology would identify the \(\mathbf F_r\)-homology objects of the condensed
homotopy types of \(\PP^1_{\mathbf C}\) and \(\PP^1_{\mathbf Q}\).

The first homology of the complex-side object has a quotient of cardinality
\(2^{2^{\aleph_0}}\): it receives the quotient of a product of copies of
\(\mathbf F_r\), indexed by \(\mathbf C\), by the corresponding direct sum.
The rational-side first homology is a quotient of a free profinite group on a
countable set and has cardinality at most \(2^{\aleph_0}\).  The two homology
objects therefore cannot be isomorphic.
\end{proof}

\begin{corollary}
There is no Grothendieck six-functor formalism on
\(X\mapsto\Dnaive(X,\Zl)\) having the usual pro-etale pullback and direct
image, arbitrary proper base change, and \(p_!\simeq p_*\) for proper maps.
The failure is not repaired by restricting merely to derived
\(\ell\)-complete objects.
\end{corollary}

\begin{proof}
Such a formalism would imply the proper base-change statement contradicted by
\cref{thm:naive-proetale-pbc-fails}.  Moreover \(\mathbf F_\ell\) is already
derived \(\ell\)-complete, so the same counterexample belongs to the proposed
complete subcategory.
\end{proof}

\begin{auditbox}[title={Why the ambient category must change}]
The obstruction is not that \(\Zl\) fails to be an honest pro-etale sheaf,
nor that completion has been omitted.  It is that arbitrary pro-etale
hypersheaves retain condensed-homotopical variation not controlled by
ordinary etale geometry.  Proper pushforward does not satisfy the base-change
axiom on all of them.  Thus ``all derived-complete pro-etale sheaves'' is not
a valid Route B coefficient system.
\end{auditbox}

\section{Solid pro-etale complexes}

Let \(X\) be qcqs.  If
\(j:U=\varprojlim_\alpha U_\alpha\to X\) is pro-etale with every
\(U_\alpha\to X\) affine etale, set
\[
 j^{\blacksquare}_{\sharp}\widehat{\mathbf Z}_{P}
   :=\varprojlim_\alpha
      (j_\alpha)_\sharp\widehat{\mathbf Z}_{P}.
\]
The solidification maps compare this continuous object with the free
pro-etale sheaf represented by \(U\).

\begin{definition}[Solid ambient category]
\label{def:solid-ambient-category}
A pro-etale sheaf \(F\) of \(\widehat{\mathbf Z}_{P}\)-modules is
\emph{solid} if, for every \(j\) as above, the natural map
\[
 \operatorname{Hom}
  \bigl(j^{\blacksquare}_{\sharp}\widehat{\mathbf Z}_{P},F\bigr)
   \longrightarrow F(U)
\]
is an equivalence.  Write
\[
 \operatorname{Sh}(X,\widehat{\mathbf Z}_{P})^{\blacksquare}
  \subset
 \operatorname{Sh}(X_{\proet},\widehat{\mathbf Z}_{P})
\]
for the full subcategory of solid sheaves.  A complex is solid if all its
cohomology sheaves are solid, and
\[
 \Dsolid(X,\widehat{\mathbf Z}_{P})
  \subset
 \Dnaive(X,\widehat{\mathbf Z}_{P})
\]
denotes the resulting full stable subcategory.  If \(\Lambda\) is a solid
\(\widehat{\mathbf Z}_{P}\)-algebra, define
\[
 \Dsolid(X,\Lambda)
   :=\operatorname{Mod}_{\Lambda}
      \bigl(\Dsolid(X,\widehat{\mathbf Z}_{P})\bigr).
\]
\end{definition}

\begin{theorem}[Solidification]
\label{thm:solidification-properties}
The inclusion
\[
 \Dsolid(X,\Lambda)\hookrightarrow\Dnaive(X,\Lambda)
\]
has a colimit-preserving left adjoint \((-)^{\blacksquare}\).  The solid
category is stable, presentable, closed symmetric monoidal, and compatible
with pullback.  It contains the pro-category generated by constructible
etale \(P\)-torsion sheaves, but it is much larger than the bounded
constructible category.
\end{theorem}

\begin{proof}
At the abelian level, solidity is precisely the reflective localization at
the solidification maps above.  The localized subcategory is closed under
extensions, all limits, and all colimits; compact generators are obtained
from cofiltered limits of constructible etale \(P\)-torsion sheaves.  Passing
to complexes gives the stated presentable stable category, and module objects
give general solid coefficients.  This is
\cite[Theorem~1.2.8 and Corollary~1.3.2]
{RuimyTubachWolfSolid}.
\end{proof}

\begin{warning}
``Solid'' is neither synonymous with ``constructible'' nor with ``derived
\(\ell\)-complete''.  For example, the topological field \(\Ql\) is a solid
\(\Zl\)-algebra.  Conversely, merely imposing derived completion does not
remove the counterexample of \cref{thm:naive-proetale-pbc-fails}.
\end{warning}

\begin{proetalebox}[title={What Route B gains}]
The coefficient object \(\Zl\), completed tensor product, scalar extension
to \(\Ql\), and all ambient adjunctions live in one presentable category.
No Artin--Rees quotient and no compatible-system object occur in the
definition.  Solidification is the additional continuity condition required
by proper geometry.
\end{proetalebox}

\section{The ambient six operations}

\begin{theorem}[Solid proper and smooth base change]
\label{thm:solid-ambient-base-change}
Consider a cartesian square of qcqs schemes
\begin{equation}
\label{eq:solid-base-change-square}
\begin{tikzcd}
X' \arrow[r,"g"] \arrow[d,"q"'] & X \arrow[d,"p"] \\
Y' \arrow[r,"f"'] & Y.
\end{tikzcd}
\end{equation}
For a solid coefficient algebra \(\Lambda\), the canonical map
\[
 f^*p_*K\longrightarrow q_*g^*K
\]
is an equivalence in either of the following cases:
\begin{enumerate}
\item \(p\) is proper;
\item \(f\) is smooth and all primes in \(P\) are invertible on \(Y\).
\end{enumerate}
\end{theorem}

\begin{proof}
We give the published proof because it answers exactly whether finite
coefficients are still needed.  The forgetful functor to
\(\widehat{\mathbf Z}_{P}\)-coefficients is conservative and commutes with
pullback and direct image, so one first reduces to that coefficient ring.
The relevant functors commute with limits and the solid standard
\(t\)-structure is left complete.  Postnikov truncation therefore reduces
to bounded-below objects.  Uniform cohomological-dimension bounds and
commutation with filtered colimits of uniformly bounded-below objects reduce
further to bounded objects, then devissage reduces to the heart.

The solid heart is generated under filtered colimits by torsion
constructible etale sheaves.  On those generators the assertion is the
ordinary torsion proper or smooth base-change theorem.  Colimit preservation
extends the equivalence to every solid object.  This is
\cite[Proposition~1.3.6]{RuimyTubachWolfSolid}.
\end{proof}

\begin{foundation}[title={Finite coefficients: definition versus proof}]
The preceding theorem is an ambient theorem: it is not defined as the limit
of the maps for \(\mathbf Z/\ell^n\).  Nevertheless, its known proof uses
torsion constructible generators and classical torsion base change.  This
cannot disappear logically: substituting \(K=\mathbf Z/\ell^n\) into the
ambient theorem recovers finite-coefficient base change.  Route B removes
the layer-by-layer \emph{construction}; it does not remove the mathematical
content of the finite theorem.
\end{foundation}

The full ambient formalism is obtained by solid rigidity.

\begin{theorem}[Ambient solid six-functor formalism]
\label{thm:ambient-solid-six-functors}
Assume all primes in \(P\) are invertible.  The assignment
\[
 X\longmapsto\Dsolid(X,\Lambda)
\]
on qcqs schemes has the following structure.
\begin{enumerate}
\item Every morphism has an adjunction \(f^*\dashv f_*\).
\item Every finitely presented morphism between qcqs schemes has an
      adjunction \(f_!\dashv f^!\).
\item If \(f\) is proper and finitely presented, then \(f_!\simeq f_*\).
\item Exceptional direct image satisfies arbitrary base change and the
      projection formula
      \[
       f_!(K\otimes f^*L)\simeq f_!K\otimes L.
      \]
\item Each fiber is closed symmetric monoidal, with tensor product and
      internal Hom.
\item Open--closed decompositions give both localization triangles.
\item The operations assemble into a lax symmetric monoidal functor on a
      category of correspondences, so all composition and base-change
      compatibilities are coherent.
\item If \(f\) is smooth, then
      \[
       f^!\simeq\Sigma_{\Omega_f}f^*.
      \]
      For a smooth map of pure relative dimension \(d\), the etale
      orientation identifies this with \(f^*(d)[2d]\).
\end{enumerate}
\end{theorem}

\begin{proof}
Ruimy--Tubach--Wolf first construct the coefficient system of solid
pro-etale motives.  Homotopy invariance, Nisnevich descent, localization,
smooth ambidexterity, and Tate stability make it a motivic coefficient
system.  The abstract correspondence theorem then supplies its six
operations; in particular, exceptional direct image is a colimit-preserving
functor and hence has the right adjoint \(f^!\).

The solid rigidity functor
\[
 \rho^{\blacksquare}_{\sharp}:
 \Dsolid(X,\Lambda)\longrightarrow
 \operatorname{DM}^{\mathrm{eff}}(X,\Lambda)^{\blacksquare}
\]
is an equivalence when all primes of \(P\) are invertible.  Full faithfulness
is proved from the solid generators, torsion rigidity, and compactness of the
Tate comparison object; essential surjectivity follows from the corresponding
motivic generators.  Transporting the motivic formalism across this
equivalence gives the theorem.  See
\cite[Theorems~2.5.4 and~3.3.7]{RuimyTubachWolfSolid}.
\end{proof}

\begin{warning}[Preprint status]
The solid rigidity theorem used here is a 2026 preprint.  It is presently the
most direct published Route B model on the pro-etale site and is especially
useful for a formalization blueprint, but its preprint status must remain
visible in the dependency ledger.  Scholze's profinitely complete
etale-valued formalism \cite[Theorems~7.15--7.16]{ScholzeSixFunctors}
provides a closely related established ambient alternative.
\end{warning}

\subsection{Compactification inside the ambient category}

Let \(f:X\to Y\) be separated and finitely presented.  Choose a Nagata
compactification
\[
 X\xrightarrow{j}\overline X\xrightarrow{p}Y
\]
with \(j\) a qcqs open immersion and \(p\) proper.

\begin{definition}[Exceptional direct image]
\label{def:proetale-fshriek}
In the ambient solid category set
\begin{equation}
\label{eq:fshriek-compactification}
 f_!K:=p_*j_!K.
\end{equation}
\end{definition}

\begin{theorem}[Independence and coherence]
\label{thm:fshriek-independent}
Definition~\eqref{eq:fshriek-compactification} is independent, through a
contractible space of coherent equivalences, of the compactification.  It
agrees with the functor in
\cref{thm:ambient-solid-six-functors}, is functorial in \(f\), satisfies
arbitrary base change, and satisfies the projection formula.
\end{theorem}

\begin{proof}
Two compactifications are dominated by the closure of the diagonal image in
their fiber product.  Proper base change compares the two functors
\(p_*j_!\).  Rather than checking an unbounded list of cocycles by hand, the
correspondence-level construction in
\cite[Theorem~2.5.4]{RuimyTubachWolfSolid} simultaneously supplies all higher
compatibilities.  This is the categorical content behind independence of
compactification.
\end{proof}

\begin{warning}
For a general \(f\), this \(f_!\) is not obtained by deriving an abelian-level
pro-etale functor.  It is a derived geometric operation produced by open
extension, proper pushforward, and coherent compactification.
\end{warning}

\section{The constructible subcategory inside the ambient theory}

We now specialize to \(P=\{\ell\}\), and let
\(\Lambda=\Zl\) or \(\mathcal O_E\).  Recall the intrinsic definition from
\cref{chap:proetale-coefficients}: lisse means tensor-dualizable, and
constructible means lisse after restriction to each member of a finite
constructible stratification.

\begin{theorem}[Constructibles are solid]
\label{thm:constructibles-embed-solid}
For every qcqs \(X\), there are full inclusions
\[
 \Dcons(X,\Zl)\hookrightarrow\Dsolid(X,\Zl),
 \qquad
 \Dcons(X,\Ql)\hookrightarrow\Dsolid(X,\Ql).
\]
The same holds over \(\mathcal O_E\) and \(E\) after scalar extension.
\end{theorem}

\begin{proof}
If \(K\) is lisse, dualizability shows that tensoring with \(K\) commutes
with limits.  Writing \(\Zl=\varprojlim_n\mathbf Z/\ell^n\), one obtains
\(K\simeq\varprojlim_nK/\ell^n\); the reductions are bounded torsion etale
complexes and hence solid.  Solidity is local and stable under the
open--closed gluing used by a finite stratification, proving the integral
claim.  A constructible \(\Ql\)-complex has a constructible \(\Zl\)-lattice,
so the rational claim follows by scalar extension.  This is
\cite[Proposition~1.3.11 and Corollary~1.3.12]
{RuimyTubachWolfSolid}.
\end{proof}

The full solid category has its own compact generators, and they must not be
confused with constructible complexes.  The correct compact recognition
takes place in the ind-constructible intermediate category.

\begin{definition}[Ind-constructible objects]
\label{def:indconstructible}
Let \(\Dindcons(X,\Lambda)\) be the full subcategory of pro-etale complexes
which are filtered colimits of objects of \(\Dcons(X,\Lambda)\).  Since the
solid category is closed under colimits and contains constructibles,
\[
 \Dcons(X,\Lambda)
  \subset\Dindcons(X,\Lambda)
  \subset\Dsolid(X,\Lambda).
\]
\end{definition}

\begin{theorem}[Compact recognition]
\label{thm:indconstructible-compact-recognition}
Suppose \(X\) has finite \(\Lambda\)-cohomological dimension in the sense of
Hemo--Richarz--Scholbach.  Then
\[
 \Dindcons(X,\Lambda)^\omega=\Dcons(X,\Lambda),
 \qquad
 \operatorname{Ind}\bigl(\Dcons(X,\Lambda)\bigr)
   \simeq\Dindcons(X,\Lambda).
\]
Finite-type schemes over finite or separably closed fields satisfy this
hypothesis for \(\Zl,\mathcal O_E,E\), and \(\Ql\).
\end{theorem}

\begin{proof}
For a constructible \(M\), the functor
\(\R\underline{\operatorname{Hom}}(M,-)\), followed by global sections,
commutes with filtered colimits of constructible terms.  The point is a
uniform cohomological-dimension bound; it removes the uniform lower-bound
hypothesis otherwise needed for direct image.  Thus every constructible
object is compact in \(\Dindcons\).

Conversely, write a compact \(M\in\Dindcons\) as a filtered colimit
\(M\simeq\varinjlim_iM_i\) of constructible objects.  Compactness makes
\(\id_M\) factor through some \(M_i\), so \(M\) is a retract of \(M_i\).
The constructible category is idempotent complete; hence \(M\) is
constructible.  The universal property of Ind-completion now gives the
second equivalence.  This is
\cite[Proposition~8.2 and Corollary~8.3]{HRSConstructible}.
\end{proof}

\begin{warning}
The theorem says
\(\Dcons=\Dindcons^\omega\), not
\(\Dcons=(\Dsolid)^\omega\).  The latter assertion is neither used nor made.
Solid compact generators include pro-torsion objects of a different nature.
\end{warning}

\section{The formal half of constructible descent}

Three operations restrict without a geometric finiteness theorem.

\begin{proposition}[Formal preservation]
\label{prop:formal-constructible-preservation}
Let \(f:X\to Y\) be a morphism of noetherian schemes and let
\(j:U\hookrightarrow X\) be a quasi-compact open immersion.
\begin{enumerate}
\item \(f^*\) preserves constructible complexes.
\item Completed tensor and internal Hom preserve pairs of constructible
      complexes.
\item \(j_!\) preserves constructible complexes.
\end{enumerate}
\end{proposition}

\begin{proof}
Pullback is strong symmetric monoidal and therefore preserves dualizable
objects.  Pulling back a finite stratification proves the first assertion.
For the second, choose a common refinement of the two stratifications.  On
each stratum \(K\) and \(L\) are dualizable, so
\[
 \R\underline{\operatorname{Hom}}(K,L)
   \simeq K^\vee\widehat\otimes L,
\]
which is again dualizable.  For the last assertion, adjoin the closed
complement to a finite stratification of \(U\).  The restriction of \(j_!K\)
to a stratum in \(U\) is the corresponding restriction of \(K\), while its
restriction to the complement is zero.  Equivalently, use
\cite[Corollary~4.6]{HRSConstructible}.
\end{proof}

\subsection{Completed pullback is now a restriction theorem}

\begin{definition}[Completed pullback]
\label{def:completed-pullback}
For a derived-complete constructible object \(K\), write
\begin{equation}
\label{eq:completed-pullback}
 \Lhat f^*K:=\bigl(\mathbf Lf^*K\bigr)^\wedge.
\end{equation}
\end{definition}

\begin{proposition}[Comparison with ambient pullback]
\label{prop:completed-pullback-adjunction}
Under the inclusion of
\cref{thm:constructibles-embed-solid}, the restriction of the ambient solid
pullback is canonically equivalent to \(\Lhat f^*\).  It preserves
constructibility and is left adjoint to the restricted \(Rf_*\) once the
latter is shown to preserve constructibility.
\end{proposition}

\begin{proof}
The ambient pullback is obtained by pullback followed by the relevant
coefficient localization; on a complete constructible object this is exactly
derived completion.  The first part of
\cref{prop:formal-constructible-preservation} gives preservation.  The
ambient adjunction restricts after the direct-image preservation theorem
below.  Reduction modulo \(\ell^n\) verifies the familiar comparison
\[
 (\Lhat f^*K)/\ell^n\simeq
 \mathbf Lf^*(K/\ell^n),
\]
but this identity is a compatibility theorem, not the definition of the
ambient operation.
\end{proof}

\section{The geometric half of constructible descent}

Ambient existence alone cannot force a functor to preserve a small
geometric subcategory.  The remaining restriction uses exactly the
Grothendieck--Gabber finiteness and constructible-duality content.

\begin{theorem}[Geometric restriction core]
\label{thm:torsion-finiteness-input}
Let \(k\) be finite or separably closed, let \(X,Y\) be finite-type
\(k\)-schemes, and let coefficients be \(\Zl\) or \(\mathcal O_E\), with
\(\ell\ne\operatorname{char}k\).
\begin{enumerate}
\item If \(p:X\to Y\) is proper and \(K\in\Dcons(X,\Lambda)\), then on each
      irreducible component of \(Y\) there is a dense open \(V\) on which
      \((p_*K)|_V\) is lisse and perfect.  Its formation is compatible with
      geometric base change on \(V\), and its cohomological amplitude is
      bounded in terms of \(p\) and the amplitude of \(K\).
\item If \(j:U\hookrightarrow X\) is a quasi-compact open immersion, then
      \(j_*\) preserves constructibility.
\item For \(a_X:X\to\Spec k\), the object
      \(\omega_X:=a_X^!\Lambda\) is constructible, the functor
      \[
       \mathbb D_XK:=
       \R\underline{\operatorname{Hom}}(K,\omega_X)
      \]
      preserves constructibles, and
      \(K\to\mathbb D_X\mathbb D_XK\) is an equivalence for constructible
      \(K\).
\end{enumerate}
\end{theorem}

\begin{proof}[Proof capsule and exact foundational status]
The first two assertions are the generic-local-acyclicity and boundary
pieces of the Grothendieck--Gabber constructibility theorem.  The proof is a
simultaneous noetherian induction on the support.  After shrinking an
irreducible base, generic local acyclicity makes the proper pushforward
locally constant and base-change compatible.  Cohomological-dimension bounds
give boundedness; finite generation on geometric fibers gives perfectness.
On the closed complement, the localization triangle and the induction
hypothesis control the boundary.  The same induction, applied to a
compactification of an open immersion, gives boundary finiteness.

For duality, compactify the structural map, use proper duality and the two
localization triangles, and induct on a regular stratification.  On a smooth
stratum of dimension \(d\), purity gives
\(\omega\simeq\Lambda(d)[2d]\); biduality is then perfect-module biduality.
The localization triangles extend it across the complement.

These are deep geometric theorems, assembled for complete coefficients in
\cite[Section~6.7]{BhattScholzeProEtale} from the classical finiteness,
local acyclicity, duality, and purity theorems.  Their currently available
proof uses finite torsion coefficients and uniform bounds.  The present notes
do not relabel that proof as ``formal solidification''.  A proof avoiding
every torsion theorem would be a new proof of the torsion special cases as
well.
\end{proof}

The global preservation theorem is now a short Route B argument in the
ambient category.

\begin{proposition}[Proper direct image preserves constructibility]
\label{prop:route-b-proper-constructible}
If \(p:X\to Y\) is proper between finite-type \(k\)-schemes, then
\[
 p_*:\Dcons(X,\Lambda)\longrightarrow\Dcons(Y,\Lambda).
\]
\end{proposition}

\begin{proof}
Proceed by noetherian induction on closed subsets of \(Y\).  For
\(K\in\Dcons(X,\Lambda)\), choose, by
\cref{thm:torsion-finiteness-input}, a dense open immersion
\(v:V\hookrightarrow Y\) such that \(v^*p_*K\) is lisse.  Let
\(i:Z\hookrightarrow Y\) be the reduced closed complement and form
\[
\begin{tikzcd}
X_Z \arrow[r,"i'"] \arrow[d,"p_Z"'] & X \arrow[d,"p"] \\
Z \arrow[r,"i"'] & Y.
\end{tikzcd}
\]
Ambient solid proper base change gives
\[
 i^*p_*K\simeq p_{Z*}i'^*K.
\]
The right side is constructible by the induction hypothesis.  Thus \(p_*K\)
is lisse on \(V\) and constructible on \(Z\); the finite-stratification
gluing criterion makes it constructible on \(Y\).
\end{proof}

\begin{corollary}[Integral direct and exceptional images]
\label{cor:integral-direct-image}
For every separated finite-type morphism \(f:X\to Y\),
\[
 f_!:\Dcons(X,\Lambda)\longrightarrow\Dcons(Y,\Lambda).
\]
Moreover, after constructible Verdier duality is applied below, the ambient
\(f_*\) also restricts to constructibles.
\end{corollary}

\begin{proof}
Choose \(f=pj\) with \(j\) open and \(p\) proper.  By
\cref{prop:formal-constructible-preservation}, \(j_!K\) is constructible;
by \cref{prop:route-b-proper-constructible}, \(p_*j_!K\) is constructible.
The ambient identity \(f_!\simeq p_*j_!\) proves the claim.
\end{proof}

\section{Verdier duality descends the remaining functors}

The ambient formalism supplies natural Verdier exchange transformations.
On reflexive constructible objects they are equivalences.

\begin{theorem}[Constructible exceptional pullback and direct image]
\label{thm:proetale-fupper-shriek}
Let \(f:X\to Y\) be separated and of finite type over \(k\).  For
constructible \(K\) there are natural ambient equivalences
\begin{equation}
\label{eq:six-duality-identities}
 f^!K\simeq\mathbb D_Xf^*\mathbb D_YK,
 \qquad
 f_*K\simeq\mathbb D_Yf_!\mathbb D_XK.
\end{equation}
Consequently both \(f^!\) and \(f_*\) preserve constructibility.
\end{theorem}

\begin{proof}
Since \(a_X=a_Yf\), transitivity gives
\(\omega_X\simeq f^!\omega_Y\).  For arbitrary ambient \(M\) and constructible \(K\), the first
Verdier exchange is verified by the chain
\begin{align*}
 \operatorname{Map}_X
  \bigl(M,\mathbb D_Xf^*\mathbb D_YK\bigr)
 &\simeq
 \operatorname{Map}_X
  \bigl(M\widehat\otimes f^*\mathbb D_YK,\omega_X\bigr)\\
 &\simeq
 \operatorname{Map}_Y
  \bigl(f_!M\widehat\otimes\mathbb D_YK,\omega_Y\bigr)\\
 &\simeq
 \operatorname{Map}_Y
  \bigl(f_!M,\mathbb D_Y\mathbb D_YK\bigr)\\
 &\simeq \operatorname{Map}_Y(f_!M,K)\\
 &\simeq \operatorname{Map}_X(M,f^!K).
\end{align*}
The second line uses the projection formula and \(f_!\dashv f^!\), and the
fourth uses constructible biduality.  This calculation is the mate
verification of the ambient Verdier exchange map, hence proves the first
equivalence.  Applying the same calculation to the other adjunction, or
taking Verdier duals, gives the second.

Every functor on the right sides of
\eqref{eq:six-duality-identities} preserves constructibility by
\cref{prop:formal-constructible-preservation,%
thm:torsion-finiteness-input,cor:integral-direct-image}.  Therefore the two
left sides are constructible.
\end{proof}

\begin{proposition}[Tensor, internal Hom, and projection]
\label{prop:proetale-tensor-hom}
For constructible \(K,L\), completed tensor product and internal Hom are
constructible, and for constructible \(M\) on \(Y\) there is an equivalence
\begin{equation}
\label{eq:fshriek-projection}
 f_!(K\widehat\otimes\Lhat f^*M)
   \simeq f_!K\widehat\otimes M.
\end{equation}
\end{proposition}

\begin{proof}
Preservation was proved in
\cref{prop:formal-constructible-preservation}.  The projection formula is an
ambient equivalence from \cref{thm:ambient-solid-six-functors}; since both
sides are constructible, it restricts without a new proof.
\end{proof}

\begin{theorem}[Purity inputs]
\label{thm:proetale-purity-inputs}
Let coefficients be as above.
\begin{enumerate}
\item If \(f:X\to Y\) is smooth of pure relative dimension \(d\), the
      etale orientation gives
      \[
       f^!K\simeq\Lhat f^*K(d)[2d].
      \]
\item If \(i:Z\hookrightarrow X\) is a regular immersion of pure codimension
      \(c\) between regular noetherian schemes, then
      \[
       i^!\Lambda\simeq\Lambda(-c)[-2c].
      \]
\end{enumerate}
\end{theorem}

\begin{proof}
The first assertion is the oriented form of smooth ambidexterity in the
ambient formalism.  The second is Gabber absolute purity.  Solidification
does not prove absolute purity: its proof uses deformation to the normal
cone, the zero-section calculation, and the corresponding torsion purity
theorem.  See \cite{ILOTravauxGabber}.
\end{proof}

\section{The restricted coherent package}

Return to a cartesian square
\begin{equation}
\label{eq:six-base-change-square}
\begin{tikzcd}
X' \arrow[r,"g'"] \arrow[d,"f'"'] & X \arrow[d,"f"] \\
Y' \arrow[r,"g"'] & Y.
\end{tikzcd}
\end{equation}

\begin{theorem}[Constructible base change]
\label{thm:proetale-base-change}
For constructible complexes:
\begin{enumerate}
\item if \(f\) is proper, then
\[
 \Lhat g^*f_*K\xrightarrow{\sim}f'_*\Lhat{g'}^*K;
\]
\item for separated finite-type \(f\), exceptional direct image satisfies
\[
 \Lhat g^*f_!K\xrightarrow{\sim}f'_!\Lhat{g'}^*K;
\]
\item the smooth base-change and projection-formula equivalences also
restrict to constructible objects.
\end{enumerate}
\end{theorem}

\begin{proof}
These are the ambient equivalences of
\cref{thm:solid-ambient-base-change,%
thm:ambient-solid-six-functors}.  The preceding sections prove that every
object appearing on either side is constructible.  Fullness of the
constructible inclusion therefore restricts the equivalences.  No separate
coherence check modulo \(\ell^n\) is necessary.
\end{proof}

\begin{theorem}[Route B constructible six operations]
\label{thm:proetale-six-functor-package}
For finite-type schemes over a finite field and separated finite-type
morphisms, the ambient solid operations restrict to
\[
 X\longmapsto\Dcons(X,\Lambda)
\]
and give
\[
 \Lhat f^*,\quad f_*,\quad f_!,\quad f^!,\quad
 -\widehat\otimes-,\quad
 \R\underline{\operatorname{Hom}}(-,-).
\]
They carry both adjunctions, composition equivalences, proper and
exceptional base change, projection formulas, localization triangles,
smooth and absolute purity where applicable, and constructible Verdier
duality.  Every higher coherence diagram is inherited from the ambient
correspondence formalism.
\end{theorem}

\begin{proof}
Pullback, tensor, internal Hom, and open extension by zero preserve
constructibility by
\cref{prop:formal-constructible-preservation}.  Proper direct image and
exceptional direct image preserve it by
\cref{prop:route-b-proper-constructible,cor:integral-direct-image}.
Exceptional pullback and ordinary direct image preserve it by
\cref{thm:proetale-fupper-shriek}.  Hence all ambient functors restrict.

The constructible categories are full subcategories.  Units, counits,
exchange maps, and projection maps are therefore the restrictions of the
ambient natural transformations.  Every equation between their composites
already holds in the ambient correspondence category, so all coherence
diagrams continue to commute after restriction.
\end{proof}

\subsection{Rational coefficients}

For a finite extension \(E/\Ql\), the safe constructible relation is
\[
 \Dcons(X,\mathcal O_E)[1/\ell]
   \simeq\Dcons(X,E).
\]
Define
\[
 \Dcons(X,\Qbarl)
   :=\varinjlim_{E/\Ql\ \mathrm{finite}}\Dcons(X,E).
\]
Linearity and finite coefficient-field descent extend the restricted six
operations to \(E\) and \(\Qbarl\).  We do not identify the rationalization
of an arbitrary derived-complete naive category with all
\(\Ql\)-sheaves; the solid topological coefficient algebra is the ambient
rational object.

\begin{powerbox}[title={Why this architecture helps the perverse proof}]
The perverse \(t\)-structure is imposed only after the ambient operations
have been constructed and after their constructible preservation has been
proved.  Consequently recollement, proper pushforward, exceptional
pullback, duality, Fourier kernels, and all base-change maps used later are
restrictions of one coherent formalism.  Perverse truncation need not carry
compactification choices or compatible \(\ell^n\)-systems in its data.
\end{powerbox}

\begin{formalization}
The reusable implementation should expose four nested layers:
\[
\begin{tikzcd}[row sep=small]
\Dnaive(X,\Lambda)
 \arrow[d,"(-)^{\blacksquare}"]\\
\Dsolid(X,\Lambda)
 \arrow[d,phantom,"\supset"]\\
\Dindcons(X,\Lambda)
 \arrow[d,phantom,"\supset"]\\
\Dcons(X,\Lambda)=\Dindcons(X,\Lambda)^\omega.
\end{tikzcd}
\]
The ambient correspondence functor and its coherence belong to the solid
layer.  Constructibility is a predicate with preservation theorems, not a
field in the ambient functor record.  The code must not assert
\(\Dcons=(\Dsolid)^\omega\), must not identify solidity with completion, and
must retain the preprint/foundational status of solid rigidity and geometric
finiteness separately.
\end{formalization}

\begin{auditbox}[title={Final Route B dependency audit}]
The exact dependency chain is
\[
\begin{tikzcd}[column sep=small,row sep=large]
\substack{\text{solidification of the}\\
          \text{naive pro-etale category}}
 \arrow[r] &
\substack{\text{solid pro-etale motives}\\
          \text{and rigidity}}
 \arrow[r] &
\substack{\text{ambient coherent}\\
          \text{six operations}}\\
\substack{\text{torsion proper base change}\\
          \text{and torsion rigidity}}
 \arrow[ur] \arrow[r] &
\substack{\text{Grothendieck--Gabber}\\
          \text{finiteness and duality}}
 \arrow[r] &
\substack{\text{restriction to}\\
          \Dcons}
 \arrow[u].
\end{tikzcd}
\]
Thus Route B genuinely constructs the operations before constructibility and
does not define them level by level.  It does \emph{not} prove that classical
finite-coefficient geometry was unnecessary.  In the current proof that
geometry enters twice: to establish ambient solid base change on generators,
and to prove that direct images and duality preserve constructibles.
\end{auditbox}
